\scanpage{12}

(en n'utilisant \unsure{pas|que} le fait que $\unsure{A_a|Aa}$ est une jacobienne) nous dit que \deleted{cette courbe} \inserted{$C_a$} a \unsure{ttes|bien} ses composantes de multiplicité $1$, et que $J^0_{C_a/k(a)}$ est isomorphe au produit des jacobiennes des composantes de $C_a$.]
% slanted
\margin{On n'a pas besoin de Hoyt, mais seulement de Matsusaka}

Mais un résultat de Matsusaka \inserted{(``géométrique'')} nous apprend que si on peut trouver sur $A$ un diviseur \inserted{\uline{positif}} \uline{irréductible} $\Theta_a$ tel que \deleted{$\deg(\Theta^n/n!)=1$, et que} $\frac{1}{n!}\deg\Theta_{a}^n=1$, \deleted{$\Theta^{n-1}/(n-1)!\ \sim$} et que $\Theta_a^{n-1}\equiv(n-1)!\,Z$ (équivalence numérique) où $Z$ est un cycle positif, alors $Z$ est \deleted{\unsure{ir}} une courbe non singulière irréductible, [et de plus -- mais nous n'en avons pas besoin -- que $A$ est isomorphe à sa jacobienne \deleted{\unsure{]}} par l'immersion $C\mathbin{\unsure{-|\to}}A$, enfin $\Theta$ définit la polarisation can. de la jacobienne].
% slanted
\margin{la première condition signifie que $\Theta$ définit une polarisation \uline{\unsure{unité}}}
\deleted{Or} \deleted{\unsure{Or} l'hypothèse \unsure{sur}} Considérons le diviseur $\Theta_\eta$ correspondant à $\unsure{J_\eta|\mathscr J_\eta}$, et son adhérence $\Theta$, d'où un $\Theta_a$. On \uline{sait} que $\deg(\Theta_\eta)^n/n!=1$ et $\Theta_\eta^{n-1}\equiv\deleted{\unsure{C_\eta}}(n-1)!\,C_\eta$, d'où $\deg(\Theta_a^n)/n!\unsure{\}|)}=1$ et $\Theta_a^{n-1}\equiv(n-1\unsure{\}|)}!\,C_a$. \deleted{\unsure{Comme}} \deleted{Si $C_a$ n'était irréductible, $\Theta_a$ ne serait pas}

% archive sheet number: 11
